\section{Configuration selection and development history}
\label{sec:selection}
\subsection{Time partitions and evaluation exposure}
Ausgrid development uses seven-day blocks from July 2010 to June 2012. Every eighth block, starting at index four, is reserved for validation; the July 2012--June 2013 year is used for evaluation. Initial results on that evaluation year were inspected before FedGrid was revised. The later numerical selector therefore uses validation scores within a development process already exposed to evaluation outcomes. New learning seeds on the same profiles do not restore an untouched evaluation year.

ACT NextGen uses 52 weeks beginning on 1 January 2018. Every fourth week, starting at index three, is assigned to evaluation, giving 13 weeks. Every eighth week, starting at index one, is assigned to validation, giving seven weeks. The remaining 32 weeks form the training partition. ACT experiments include retraining and transfer without retraining, while retaining the feeder, tariff and battery-control assumptions. The donor audit in the appendix additionally documents imputation across these partition boundaries.

\subsection{Validation criterion}
For each configuration and seed, selection averages the final four validation assessments. Each assessment sums household-average cost and the curtailment penalty across the validation weeks. Validation damage is the attacked score minus the matched no-attack score. Configurations are ranked by the largest attack-wise mean of these paired differences.

FedGrid candidates use flip, scale, noise and Normalized attacks at a nominal 30\% compromised share and seeds 0--2. The tuned TAFDRL-style comparator uses flip and Normalized only. The candidate families consequently have unequal tuning attack sets and search budgets. Their cost-plus-curtailment scores over validation weeks are distinct from the annual evaluation bill damage.

\subsection{Post-hoc rejection constraint}
The robustness-only selector initially chose a configuration applying trust separately to each household. Its worst-case mean validation damage was \$6.2 per household over the validation weeks, and its recorded honest rejection was 30.8\%, compared with 2.0\% for original FedGrid. A numerical eligibility constraint was then added: a FedGrid configuration could reject at most twice the original configuration's recorded honest share, giving a 4.0\% bound for these results. This bound was introduced after the candidate outcomes had been inspected; it was not preregistered.

The development record shows that a settings sweep for the household-trust configuration had already finished before the change. The subsequent sweep of the newly selected clipping configuration followed the revised criterion. Thus, the eligibility rule preceded the later sweep, while the earlier settings results were already available.

Under the revised criterion, filtering before clipping with a multiplier of two was selected. Its worst-case validation damage was \$14.5 and recorded honest rejection was 3.4\%. A one-setting-at-a-time sweep considered two or four solar bands and outlier thresholds of 2.0 or 3.0. The default three bands and threshold 2.5 remained selected. The tuned TAFDRL-style filter used $\lambda=2$, minimum trust 0.35 and chi-square level 0.99. Table~\ref{tab:selection_candidates} and Figure~\ref{fig:selection} report the development outcomes underlying the choice.

\subsection{Scope of the final comparison}
The final configurations were frozen after the available evaluation outcomes had been inspected. The present analysis uses that fixed archive, including unfavourable results, without further algorithm changes. Paired effects and exact tests are retrospective analyses of existing runs. The eligibility bound concerns recorded rejection, whose relationship to effective contribution is derived in Section~\ref{sec:methods} and qualified in Section~\ref{sec:participation}.
